\section{The global decomposition}
\label{sec:global}

We now assume that $G$ is semisimple, and put
$U=X\setminus\operatorname{supp}(D)$.  Throughout this section,
Frobenius and the normalized Satake isomorphism have the conventions
fixed in the introduction.  In particular, if $d_x=\deg(x)$, then
$q_x=q^{d_x}$.

Theorem~\ref{ext:discrete} gives a complex compact--Jacobson--Morozov
form at each good place and each complex realization.  Three steps
remain.  At each place we must obtain one algebraic factorization whose
remainder is compact at every complex embedding.  We must then show
that its nilpotent orbit is independent of the good place, and descend
the resulting decomposition of joint Hecke eigenspaces to $\mathbb Q$.

\subsection{The rational cuspidal spectrum}

We first record precisely the finite-dimensional space to which the
argument applies.  For a split group, the isomorphism classes of bundles
with full level $D$ are
\[
 \Bun_{G,D}(\mathbb F_q)
   =G(F)\backslash G(\mathbb A_F)/K_D.
\]
In general the bundle description has a disjoint union indexed by
$\ker^1(F,G)$; this kernel is trivial for split $G$
\cite[Section~2.1, Equation~(2.1), and Remark~2.2]{VLafforgue}.
Thus this identification does not assert the stronger, generally false
statement that every $G$-torsor over $F$ is trivial.
The finite-level cuspidal space is finite dimensional
\cite[Section~1.1]{VLafforgue}.  The finite subscheme called $N$ there
can be precisely $D$, with all its multiplicities: its subgroup
$K_N$ is defined by the same kernel as $K_D$.  Moreover, the central
lattice used there can be trivial here, since the center of a semisimple
group has compact adelic points.  These statements therefore concern
the entire space $C_{D,\mathbb Q}$ in Theorem~\ref{main:decomposition}.

Write $C_{D,\mathbb C}=\mathbb C\otimes_{\mathbb Q}C_{D,\mathbb Q}$,
and let
\[
 \mathbb T_{D,\mathbb Q}
 \subset\operatorname{End}_{\mathbb Q}(C_{D,\mathbb Q})
\]
be the unital algebra generated by the good-place spherical Hecke
algebras $A_x$, for $x\in |U|$.

\begin{lemma}\label{glob:hecke-spectrum}
The algebra $\mathbb T_{D,\mathbb Q}$ is finite dimensional and reduced.
After extending scalars to $\overline{\mathbb Q}$, its action has a
decomposition
\begin{equation}\label{glob:joint-spectrum}
 \overline{\mathbb Q}\otimes_{\mathbb Q}C_{D,\mathbb Q}
   =\bigoplus_{\chi\in\Sigma_D} C_\chi
\end{equation}
into joint eigenspaces, where $\Sigma_D$ is finite.  All good-place
Satake classes of these characters are algebraic.
\end{lemma}

\begin{proof}
Finite dimensionality of $C_{D,\mathbb Q}$ implies that of its Hecke
image algebra.  On the complex cusp space use the positive automorphic
$L^2$ inner product.  For a spherical function $a$, the adjoint of
convolution by $a$ is convolution by
$a^*(g)=\overline{a(g^{-1})}$.  This function is again spherical and
its operator preserves the cusp space.  The spherical operators at
one place commute, and operators at different places commute.
Consequently all their restrictions are commuting normal operators.
Simultaneous diagonalization follows in the finite-dimensional space,
and implies that the Hecke image algebra is reduced.  A finite reduced
algebra over $\mathbb Q$ is a product of number fields, so it splits
over $\overline{\mathbb Q}$ and gives
Equation~\eqref{glob:joint-spectrum}.  The normalized Satake
isomorphism is defined over $\mathbb Q(\sqrt{q_x})$; hence an algebraic
Hecke character determines an algebraic point of the dual adjoint
quotient, and therefore an algebraic semisimple conjugacy class.
\end{proof}

\subsection{A single triple at every complex embedding}

The definition of $R_{r,\mathcal O}$ requires one algebraic triple and
an algebraic remainder compact at every complex embedding.  The local
theorem supplies complex triples separately at those embeddings.  We
compare their neutral cocharacters in a fixed torus: the exact root
equations will force them to agree, after which polynomial equations
produce a triple over $\overline{\mathbb Q}$.

For a split maximal torus $T$ of $\Gd$ and a semisimple
$t\in T(\overline{\mathbb Q})$, define its exponent at an embedding
$\iota:\overline{\mathbb Q}\hookrightarrow\mathbb C$, relative to a
prime power $r>1$, by
\begin{equation}\label{glob:exponent}
 \langle a,\lambda_\iota(t)\rangle
   =\log_r|\iota(a(t))|,
 \qquad a\in X^*(T).
\end{equation}
Thus $\lambda_\iota(t)\in X_*(T)\otimes_{\mathbb Z}\mathbb R$.
We first keep this vector in the fixed torus, without applying a
different Weyl element at each embedding.

\begin{lemma}\label{glob:embeddings}
Suppose that, at every embedding $\iota$, the element $\iota(t)$ is
conjugate to
\[
 \phi_\iota\!\begin{pmatrix}\sqrt r&0\\0&\sqrt r^{-1}\end{pmatrix}
 c_\iota,
\]
where $\phi_\iota:\operatorname{SL}_2\to\Gd_{\mathbb C}$ is an
algebraic homomorphism and $c_\iota$ is semisimple and conjugate into
a compact subgroup of
$Z_{\Gd}(\phi_\iota(\operatorname{SL}_2))(\mathbb C)$.
Then the following hold.
\begin{enumerate}[label=\textup{(\roman*)}]
\item The vectors $\lambda_\iota(t)$ are all equal to a vector
      $\lambda$, and $\mu=2\lambda$ belongs to $X_*(T)$.
\item There is a homomorphism
      $\phi:\operatorname{SL}_2\to\Gd_{\overline{\mathbb Q}}$
      with diagonal cocharacter $\mu$ such that
      \[
        t=\mu(\sqrt r)c,
        \qquad
        c\in Z_{\Gd}(\phi(\operatorname{SL}_2))
                  (\overline{\mathbb Q}),
      \]
      and the image of $c$ is compact at every complex embedding.
\item The nilpotent orbit of the raising element of $\phi$ is unique.
      In particular, $[t]$ belongs to precisely one of the sets
      $R_{r,\mathcal O}$.
\end{enumerate}
\end{lemma}

\begin{proof}
Fix one embedding and conjugate its proposed decomposition so that
its product is $\iota(t)$ itself.  Let $A$ be the image of the diagonal
torus of $\operatorname{SL}_2$.  Its centralizer
$L=Z_{\Gd}(A)$ is a connected Levi subgroup
\cite[Theorem~17.38 and Section~25.25]{MilneAlgebraicGroups}.
The semisimple element
$c_\iota$ lies in $L$, so it lies in a maximal torus $S$ of $L$.
The central torus $A$ belongs to every maximal torus of $L$.
Therefore $S$ contains both $A$ and $\iota(t)$ and is a maximal torus
of $\Gd$.  Both $S$ and $T_{\mathbb C}$ are maximal tori of the
connected group $Z_{\Gd}(\iota(t))^\circ$; maximal-torus conjugacy
\cite[Theorem~17.10]{MilneAlgebraicGroups} therefore puts the diagonal
cocharacter in $T$ without moving $\iota(t)$.
This argument never requires the full semisimple
centralizer, or the centralizer of the triple, to be connected.

After this alignment, let $\mu_\iota\in X_*(T)$ be the diagonal
cocharacter.  The compact factor belongs to $T(\mathbb C)$, and the
closure of its cyclic subgroup is compact.  Every torus character
therefore has modulus one on it.  Equation~\eqref{glob:exponent}
gives
\begin{equation}\label{glob:integral-neutral}
             2\lambda_\iota(t)=\mu_\iota.
\end{equation}

We prove that these integral cocharacters are independent of the
embedding.  Choose a positive definite Weyl-invariant inner product
on $X_*(T)\otimes\mathbb R$, and extend it to an invariant complex
bilinear form on $\gd_{\mathbb C}$.  Such a form is obtained by
choosing a positive multiple of the Killing form on each simple
factor.  For two embeddings write $\lambda,\lambda'$ for their
exponents and $\delta=\lambda'-\lambda$.  The aligned triple at the
first embedding has neutral element
$H_\mu=d\mu(1)=2\lambda$ under the usual identification of the
real cocharacter space with the real subspace of $\operatorname{Lie}T$.
If a root component $e_\alpha$ of its raising element is nonzero,
then
\[
       \iota(\alpha(t))=r,
       \qquad \alpha(\lambda)=1.
\]
Since $\alpha(t)$ is algebraic and $r$ is rational, this is the
algebraic identity $\alpha(t)=r$.  At the second embedding it implies
$\alpha(\lambda')=1$.  The corresponding argument for the lowering
element uses the identity $\alpha(t)=r^{-1}$.  It follows that
$\delta$ centralizes both the raising and lowering elements, say
$e$ and $f$.  Invariance gives
\[
        (\delta,2\lambda)
        =(\delta,[e,f])=([\delta,e],f)=0.
\]
Applying the same argument to the aligned triple at the second
embedding gives $(\delta,2\lambda')=0$.  Subtracting proves
$(\delta,\delta)=0$, whence $\lambda'=\lambda$.
Equation~\eqref{glob:integral-neutral} proves \textup{(i)}.

The common cocharacter $\mu$ is defined over $\mathbb Q$, since $T$
is split.  Put $H_\mu=d\mu(1)$.  The equations in $e,f\in\gd$
\begin{equation}\label{glob:algebraic-triple}
 \begin{gathered}
 [H_\mu,e]=2e,\qquad [H_\mu,f]=-2f,\qquad [e,f]=H_\mu,\\
 \operatorname{Ad}(t)e=re,\qquad
 \operatorname{Ad}(t)f=r^{-1}f
 \end{gathered}
\end{equation}
are polynomial equations over $\overline{\mathbb Q}$ and have a
complex solution by the alignment above.  They consequently have a
solution over $\overline{\mathbb Q}$.  In characteristic zero, the
resulting $\mathfrak{sl}_2$ homomorphism integrates to an algebraic
homomorphism $\phi:\operatorname{SL}_2\to\Gd$.  Its diagonal
cocharacter is $\mu$, because cocharacters in characteristic zero
are determined by their differentials.  This includes the zero
triple and the trivial homomorphism.

Choose an algebraic square root of $r$ and set
$c=t\mu(\sqrt r)^{-1}\in T(\overline{\mathbb Q})$.
Equation~\eqref{glob:algebraic-triple} says that $c$ centralizes the
triple, hence its connected algebraic image
$\phi(\operatorname{SL}_2)$.  At every embedding and for every
$a\in X^*(T)$,
\[
 |\iota(a(c))|
   =r^{\langle a,\lambda\rangle}
       r^{-\langle a,\mu\rangle/2}=1.
\]
Thus the cyclic subgroup generated by $\iota(c)$ has compact closure
in $T(\mathbb C)$.  This closure is contained in the closed subgroup
$Z_{\Gd}(\iota\phi(\operatorname{SL}_2))(\mathbb C)$, and is
conjugate into a maximal compact subgroup of that reductive group
\cite[Theorems~3.6 and~3.7]{AdamsTaibiCohomology}.
This argument also applies to a disconnected triple centralizer and
proves \textup{(ii)}.

For uniqueness, first observe that any compact--Jacobson--Morozov
decomposition of $t$ has diagonal cocharacter $\mu$ after the same
torus alignment.  Decompose the Lie algebra by $H_\mu$ as
$\gd=\bigoplus_j\gd_j$.  For any triple $(e,H_\mu,f)$,
elementary $\mathfrak{sl}_2$ representation theory gives
\[
                  [\gd_0,e]=\gd_2:
\]
in every irreducible $\mathfrak{sl}_2$ summand having weight two,
raising maps weight zero onto weight two.  The orbit of $e$ under
$Z_{\Gd}(H_\mu)^\circ$ is therefore open in the irreducible vector
space $\gd_2$.  There can be only one such open orbit.  Hence the
neutral cocharacter determines the nilpotent orbit, without an
exception for very-even orbits in type~$D$.

Finally conjugate $\phi$ over $\overline{\mathbb Q}$ to the selected
representative $\phi_{\mathcal O}$ of its orbit.  The conjugated
remainder satisfies precisely the all-embeddings condition defining
$H_{\mathcal O}^c(\overline{\mathbb Q})$
\cite[Sections~3.1.1 and~3.4.1--3.4.2]{GLR}.  This proves
\textup{(iii)}.
\end{proof}

\begin{remark}\label{glob:signs}
The sets $R_{r,\mathcal O}$ are pairwise disjoint by
Lemma~\ref{glob:embeddings}.  They are Galois-stable with the rational
forms stipulated in their definition
\cite[Section~3.4.2]{GLR}.  To see the square-root issue directly,
replacing $\sqrt r$ by $-\sqrt r$ inserts
$\phi_{\mathcal O}(-I)$ into the compact factor.  This element is
central in the triple centralizer and has finite order, so preserves
compactness at every embedding.  Multiplication of $t$ by any
finite-order element of $Z(\Gd)$ likewise preserves its orbit label
and membership in $R_{r,\mathcal O}$.

For later use, let $2\rho_G$ be the sum of the positive roots of $G$,
viewed as a cocharacter of the dual torus.  In the usual square-root
presentation of normalized Satake, if $\tau$ changes the
chosen square root of $q_x$ by a sign $\epsilon_{\tau,x}$, then the
Satake class of the Galois-conjugate Hecke character is
\begin{equation}\label{glob:satake-galois}
       [t_{\chi^\tau,x}]
       =[(2\rho_G)(\epsilon_{\tau,x})\,\tau(t_{\chi,x})].
\end{equation}
Indeed, on a dominant coweight $a$ the change of the normalization
factor is $\epsilon_{\tau,x}^{\langle2\rho_G,a\rangle}$.
The element $(2\rho_G)(-1)$ is central in $\Gd$, since its pairing
with every dual root is even.  It has order at most two.  Thus
Equation~\eqref{glob:satake-galois} preserves all the assertions of
Lemma~\ref{glob:embeddings}.  If the rational Satake form absorbs
this correction, the comparison is ordinary Galois equivariance.
In either presentation of the identification in
\cite[Equation~(3.2)]{GLR}, the comparison differs at most by the
finite central factor just computed; only this conclusion is used below.
\end{remark}

\subsection{Purity on the good open curve}

We next prove that the orbit found locally is independent of the
good place. Sawin and Templier's Theorem~11.7 gives an earlier
place-independence result for unramified temperedness using
V.~Lafforgue's parameters \cite{SawinTemplier}.
Here the global purity argument compares every nilpotent-orbit label.
We will use the following precise consequence of the
purity theorem, allowing ramification at every point of $X\setminus U$.

\begin{lemma}\label{glob:twist-purity}
Let $\mathcal V$ be a lisse $\overline{\mathbb Q}_\ell$-sheaf on the
smooth geometrically connected curve $U/\mathbb F_q$, defined over a
finite extension of $\mathbb Q_\ell$.  Suppose that all its Frobenius
eigenvalues at closed points are algebraic over $\mathbb Q$.
For each complex embedding of these algebraic numbers, the multiset
of numbers
\[
       \log_{q_x}|\alpha|,
       \qquad \alpha\text{ an eigenvalue of }
                        \operatorname{Frob}_x\text{ on }\mathcal V,
\]
counted with multiplicity, is independent of $x\in|U|$.
\end{lemma}

\begin{proof}
Take the irreducible constituents $\mathcal W_j$ of a
Jordan--H\"older filtration of $\mathcal V$.  After a finite extension
of coefficients they are defined over finite extensions of
$\mathbb Q_\ell$.  At each closed point the characteristic polynomial
of $\mathcal V$ is the product of those of its constituents, so their
eigenvalues are still algebraic.  It suffices to treat one constituent
$\mathcal W$ of rank $n$.

The rank-one sheaf $\det\mathcal W$ has finite geometric monodromy
\cite[Proposition~1.3.4]{DeligneII}.  Fix a closed point $y\in|U|$
of degree $d_y$ and choose an algebraic number $b$ such that
\[
                b^{n d_y}
                   =\det(\operatorname{Frob}_y\mid\mathcal W).
\]
Normalize the degree map on the Weil group by
$\deg(\operatorname{Frob}_x)=d_x$.  The number $b$ is an
$\ell$-adic unit: the original representation has compact image,
so the eigenvalues of every group element, and in particular its
determinant, are $\ell$-adic units.  The constant-field character
$b^{\deg}$ consequently extends continuously from the Weil group
to the arithmetic fundamental group.  Set
$\mathcal W^0=\mathcal W\otimes b^{-\deg}$.

Its determinant still has finite geometric monodromy, and has value
one on $\operatorname{Frob}_y$.  A power killing its geometric
monodromy factors through $\operatorname{Gal}(\overline{\mathbb F}_q/
\mathbb F_q)$; its value on an element of degree $d_y$ is one, so this
constant-field character has finite order.  Thus
$\det\mathcal W^0$ has finite order.  This is also the determinant
normalization described in \cite[Section~1.3.6]{DeligneII}.

Laurent Lafforgue's purity theorem applies to the irreducible lisse
sheaf $\mathcal W^0$ on the smooth curve $U$: its finite-order
determinant implies that all its Frobenius eigenvalues are algebraic
and have modulus one under every complex embedding
\cite[Theorem~VII.6\textup{(i)--(ii)}]{LLafforgue}.  Therefore every
Frobenius eigenvalue $\alpha$ of $\mathcal W$ at $x$ satisfies
\[
            \log_{q_x}|\alpha|=\log_q|b|.
\]
This value is independent of $x$.  Taking the union over the fixed
list of constituents, with their ranks and multiplicities, proves
the assertion.  Neither the rank-one theorem nor the purity theorem
requires $U$ to be proper.
\end{proof}

\begin{proposition}\label{glob:place}
Fix a joint character $\chi\in\Sigma_D$ and an embedding
$\iota:\overline{\mathbb Q}\hookrightarrow\mathbb C$.
Let $\lambda_{\chi,x}$ be the dominant exponent of its normalized
Satake class at $x\in|U|$, with base $q_x$ and absolute value induced
by $\iota$.  Then $\lambda_{\chi,x}$ is independent of $x$.
\end{proposition}

\begin{proof}
Choose an embedding
$\overline{\mathbb Q}\hookrightarrow\overline{\mathbb Q}_\ell$ and
view $\chi$ as an $\ell$-adic character.  V.~Lafforgue's global
parameter theorem gives a continuous semisimple parameter
\[
     \sigma:\pi_1(U)\longrightarrow\Gd(\overline{\mathbb Q}_\ell),
\]
defined over a finite extension of $\mathbb Q_\ell$, for which
the semisimple class of $\sigma(\operatorname{Frob}_x)$ is the Satake
class of $\chi$ at every $x\in|U|$
\cite[Theorem~1.1]{VLafforgue}.  Here is why one parameter can be
chosen for this joint character.  Its nonzero joint eigenspace is
preserved by the commuting excursion operators.  Their
finite-dimensional image has a character on that eigenspace.
The parameter attached to such an excursion character has the
required good-place Satake evaluations.  We need neither uniqueness
of this lift of $\chi$ nor reducedness of the excursion algebra.

For an algebraic representation $V$ of $\Gd$, the composition
$V\circ\sigma$ is a lisse sheaf on $U$.  Its Frobenius eigenvalues
are algebraic, since the corresponding semisimple Satake classes
are algebraic.  Lemma~\ref{glob:twist-purity} shows that their
multiset of logarithmic absolute values, with base $q_x$, is
independent of $x$.

For completeness, this determines the exponent in the dual torus.
Let $V_a$ have dominant highest weight $a$.  If
$\lambda_{\chi,x}$ is dominant, every weight $a'$ of $V_a$ satisfies
\[
       \langle a',\lambda_{\chi,x}\rangle
          \le\langle a,\lambda_{\chi,x}\rangle.
\]
The largest Frobenius exponent in $V_a\circ\sigma$ is therefore
$\langle a,\lambda_{\chi,x}\rangle$, and it is independent of $x$.
The dominant weights in the character lattice of the dual maximal
torus span its real weight space: a sufficiently divisible positive
multiple of each fundamental weight belongs to this lattice.
Their pairings determine
$\lambda_{\chi,x}$.  This argument applies to every isogeny type
of the semisimple dual group.
\end{proof}

\subsection{Descent of the orbit summands}

\begin{proposition}\label{glob:descent}
For every effective divisor $D$, the rational cuspidal space has the
decomposition
\begin{equation}\label{glob:rational-decomposition}
             C_{D,\mathbb Q}
                 =\bigoplus_{\mathcal O}C_{D,\mathcal O,\mathbb Q}
\end{equation}
with the spectral-support summands specified in the introduction.
\end{proposition}

\begin{proof}
If the cusp space is zero, the assertion is immediate.  Otherwise
fix $\chi\in\Sigma_D$ and a good place $x$.  At any complex
realization, a nonzero vector of $C_\chi$ occurs in the discrete
cuspidal automorphic spectrum.  Choose a nonzero projection of such a vector to an irreducible
cuspidal constituent.  This projection retains the joint Hecke
character and is spherical at $x$.
Lemma~\ref{ext:hecke} shows that its Iwahori invariants contain the
dominant ordering of its spherical Bernstein character.
Theorem~\ref{ext:discrete} therefore puts the spherical parameter
in compact--Jacobson--Morozov form.

The same statement applies to every Galois-conjugate character:
$C_{D,\mathbb Q}$ and its Hecke action are defined over $\mathbb Q$,
so every such character occurs again in
Equation~\eqref{glob:joint-spectrum}.  By
Equation~\eqref{glob:satake-galois} and Remark~\ref{glob:signs}, every
embedding of an algebraic torus representative $t_{\chi,x}$ itself
has compact--Jacobson--Morozov form.  Lemma~\ref{glob:embeddings}
then gives a unique orbit $\mathcal O_{\chi,x}$ with
\[
                [t_{\chi,x}]\in R_{q_x,\mathcal O_{\chi,x}}.
\]
Proposition~\ref{glob:place} makes its dominant exponent independent
of $x$.  The orbit uniqueness in Lemma~\ref{glob:embeddings} thus
makes $\mathcal O_{\chi,x}$ independent of $x$ as well; denote it
by $\mathcal O_\chi$.

We have partitioned the finite set $\Sigma_D$ into the subsets
$\Sigma_{D,\mathcal O}=\{\chi:\mathcal O_\chi=\mathcal O\}$.
They are Galois-stable by Remark~\ref{glob:signs} and
Equation~\eqref{glob:satake-galois}.  Let $e_\chi$ be the primitive
idempotent of
$\overline{\mathbb Q}\otimes\mathbb T_{D,\mathbb Q}$ associated
with $\chi$.  The sum
\[
            e_{\mathcal O}
                 =\sum_{\chi\in\Sigma_{D,\mathcal O}}e_\chi
\]
is Galois-fixed, so belongs to $\mathbb T_{D,\mathbb Q}$.
These idempotents are orthogonal and sum to one.  They yield a
rational direct sum decomposition with summands
$e_{\mathcal O}C_{D,\mathbb Q}$.

It remains to identify them with the given cyclic-support spaces.
For $f\in C_{D,\mathbb Q}$ write $f=\sum_\chi f_\chi$ after scalar
extension using Equation~\eqref{glob:joint-spectrum}.
At a fixed place $x$, the character support of
$\overline{\mathbb Q}\otimes A_x f$ is exactly
\[
         \{\chi|_{A_x}:f_\chi\ne0\}.
\]
Indeed, the local eigenspace for a restricted character is the
direct sum of the joint eigenspaces with that restriction, and the
projection of $f$ to it is nonzero precisely when one of the
corresponding $f_\chi$ is nonzero.  It follows that this cyclic
support is contained in $R_{q_x,\mathcal O}$ for every good $x$ if
and only if every $\chi$ with $f_\chi\ne0$ has
$\mathcal O_\chi=\mathcal O$.  Here we use the pairwise
disjointness of the orbit sets.  The condition is therefore
equivalent to $e_{\mathcal O}f=f$, proving
$e_{\mathcal O}C_{D,\mathbb Q}=C_{D,\mathcal O,\mathbb Q}$.
This is the spectral-support definition of
\cite[Section~3.4.3]{GLR}, with the closed points restricted to $U$.
Equation~\eqref{glob:rational-decomposition} follows.
\end{proof}

\begin{proof}[Proof of Theorem~\ref{main:decomposition}]
Proposition~\ref{glob:descent} proves the rational assertion.
Tensor Equation~\eqref{glob:rational-decomposition} with
$\overline{\mathbb Q}_\ell$ over $\mathbb Q$.  The resulting
summands are exactly $C_{D,\mathcal O,\ell}$ by their definition.
Scalar extension is exact, so their natural addition map to
$C_{D,\ell}$ is an isomorphism.
\end{proof}

\subsection{Generic cuspidal representations}

\begin{proof}[Proof of Corollary~\ref{main:generic-ramanujan}]
Choose the split model of $G$ over the constant field with the same
root datum. Its base change to $F$ is isomorphic to the given group.
Let $S$ be the finite set of places at which $\pi$ is ramified.
For a split adjoint group all hyperspecial maximal compact subgroups
at a place are conjugate, so unramifiedness is equivalent to having
$G(\mathcal O_x)$-fixed vectors for this model. Smoothness of the local
representations and their restricted tensor product allow an effective
divisor $D$ supported on $S$ with $\pi^{K_D}\ne0$: at each point of $S$,
take a principal congruence subgroup fixing a nonzero local vector.
The spherical Hecke algebras outside $S$ act on $\pi^{K_D}$ through the
joint character determined by the local components $\pi_x$.
This character occurs in the complex cusp space
$C_{D,\mathbb C}=\mathbb C\otimes_{\mathbb Q}C_{D,\mathbb Q}$.
By Lemma~\ref{glob:hecke-spectrum}, it and its Satake classes are
algebraic. Choose $\iota:\overline{\mathbb Q}\hookrightarrow\mathbb C$
realizing this character, and write $[t_x]$ for its algebraic normalized
Satake classes. The complex scalar extension of
Theorem~\ref{main:decomposition} and the
pairwise disjointness in Lemma~\ref{glob:embeddings} give a single orbit
$\mathcal O$ such that
\[
                 [t_x]\in R_{q_x,\mathcal O}
                 \qquad (x\notin S).
\]
We use positive real square roots in the complex normalized Satake
parameters of $\pi_x$. Their comparison with $\iota(t_x)$ has at most
the finite central correction in Remark~\ref{glob:signs}, which changes
neither the orbit nor the polar exponent.

At the place $v$, where $\pi_v$ is generic, choose an algebraic representative
$t_v\in\widehat T(\overline{\mathbb Q})$ in a maximal torus of $\Gd$.
Let $\nu_v\in X_*(\widehat T)\otimes\mathbb R$
be its polar exponent, so
$|\iota(a(t_v))|=q_v^{\langle a,\nu_v\rangle}$ for every
$a\in X^*(\widehat T)$. The torus alignment in
Lemma~\ref{glob:embeddings} identifies
\[
                  \nu_v=\tfrac12\mu_{\mathcal O},
\]
where $\mu_{\mathcal O}\in X_*(\widehat T)$ is the diagonal
cocharacter of a Jacobson--Morozov homomorphism for $\mathcal O$.
For every weight $a$ of an algebraic representation of $\Gd$,
\[
        \langle a,\nu_v\rangle
             =\tfrac12\langle a,\mu_{\mathcal O}\rangle
             \in\tfrac12\mathbb Z.
\]
Thus the complex Satake parameter of $\pi_v$ is geometric in the sense of
\cite[Definition~6.3]{CiubotaruHarris}. Since $G$ is adjoint, $\Gd$ is
simply connected; in particular, the adjoint and fundamental
representations required in their Corollary~6.32 are algebraic
representations of $\Gd$ and satisfy the displayed condition.

The cuspidal representation $\pi$ occurs in the unitary cuspidal
$L^2$ spectrum, so each local component is unitarizable. The component
$\pi_v$ is therefore spherical, generic, and unitary. By
\cite[Corollary~6.32\textup{(1)}]{CiubotaruHarris}, its polar exponent
is zero. Hence $\mu_{\mathcal O}=0$. The $\mathfrak{sl}_2$ relations
then make the raising and lowering elements zero, so
$\mathcal O=\{0\}$. For this orbit the Jacobson--Morozov homomorphism is
trivial, so every class in $R_{q_x,\{0\}}$ has compact image under
$\iota$. A spherical representation with compact normalized Satake
parameter is tempered. This proves the assertion at every $x\notin S$.

Finally, a nonzero global Whittaker functional on $\pi$ induces a
nonzero local Whittaker functional at each place. Indeed, choose a pure
tensor on which the global functional is nonzero and fix all its factors
except the chosen local one. The resulting functional on that factor
has the required local equivariance and is nonzero. A globally generic
$\pi$ is therefore generic at every place and unramified outside the
finite set $S$, so it satisfies the hypothesis of the corollary.
\end{proof}
